\section{Introduction}
\label{sec:Introduction}

A \textbf{time series} is an ordinal sequence of data points which have been sampled at a consistent rate over some period of time. \textbf{Time series classification} fits a model which maps the time series, or a subsequence of the time series, to a discrete response variable. For example, one or more sensors monitoring human physiological data could generate a time series; from the time series it might be possible to classify human behavior such as sleeping, sitting, or walking. Novel machine learning algorithms for time series classification are developed every year, and are evaluated on standard datasets maintained by academic researchers working in the field. 

To date, limited attention has been paid to the hardware acceleration of time series classification algorithms. Researchers who develop the algorithms, typically in Python, do not engage in substantive performance tuning beyond, e.g., optimizing hyperparameters. The best performing algorithms in terms of classification accuracy take days to train (on relatively limited datasets, e.g., compared to ImageNet in computer vision), and, in response, have developed more efficient algorithms which sacrifice some accuracy for computational efficiency. To the best of our knowledge, there has been no research, to date, on hardware accelerators, either FPGA or ASIC, to accelerate state-of-the-art algorithms in this domain. 

FPGAs, in particular, are a promising solution to accelerate time series classification algorithms applied to streaming data in real-time; an FPGA connected to a network can process streaming datapoints as soon as they are extracted from network packets, as opposed to CPU or GPU-based systems which would require the datapoints to be first copied to host memory as a precursor to computational analysis as imposed by modern networking protocols (e.g., TCP, Ethernet, RDMA). While a modern SmartNIC could process the time series locally without copying it to host memory, it would still need to copy the data from a discrete chip which performs packet processing to another discrete chip (e.g., a CPU or an FPGA) mounted on the NIC. Performing packet processing and data processing together within an FPGA connected to the network eliminates these sources of inefficiency. 

This paper presents an FPGA-based accelerator for the MiniRocket time series classification algorithm \cite{MiniRocket}. MiniRocket transforms an input time series into a set of features for a ridge regression or logistic regression classifier using randomly-generated (not trained) convolutional kernels followed by a pooling layer which extracts the features. MiniRocket embodies a number of properties which make it amenable to FPGA acceleration, such as fixing the length of the convolution kernels to nine scalar elements, and restricting each scalar element to be either $-1$ or $2$. \textbf{ToDo: Provide a short 1-3 sentence summary of the experimental results.}

\textbf{ToDo: Add a paragraph describing the paper organization}
